\makeabstract
{
Accretion in the Blue Ring Nebula’s Stellar Merger Remnant: TYC 2597-735-1
}
{
Valentina Guerrero Cachucho (1), Keri Hoadley, PhD (2), Jessica Li, PhD (2)
}
{
\textsuperscript{1}Department of Physics and Astronomy, Amherst College, Amherst, MA, USA\\
\textsuperscript{2}Department of Astronomy, University of Florida, Gainesville, FL, USA
}
{
Stellar mergers are a short-lived but common occurrence through the evolution of binary star systems. An opaque shell of gas and dust created during the merger obscures the remaining stellar merger remnant, making it nearly impossible to observe these objects. An ultraviolet blue ring nebula with a stellar merger remnant at its center was reported in 2020, making it the first merger system unobstructed by dust. Studying these remnants can provide new insights into the evolution of merger remnants, as well as the creation of peculiar stars like rapid rotators or blue stragglers, which are thought to come from stellar merger events.
}
{
This study sought to compare the observed spectrum of the stellar merger remnant to a model stellar atmosphere spectrum to assess the presence of excess ultraviolet continuum, a signature of active accretion occurring in a star. The observed spectrum used for this study was obtained from the Galaxy Evolution Explorer (GALEX) and spans 1300-3000 Å in the ultraviolet spectrum. The model used for comparison was selected by matching the metal abundance, surface gravity, and effective temperature parameters to the observed values. Regions with significant absorption features were masked in the observed and modeled spectra, allowing for a comparison of their continua. A chi-square goodness-of-fit test assessed whether the model spectra could predict the excess in ultraviolet emission in the observed spectrum. The goodness-of-fit test was performed for two regions: the entire observed spectrum with a binned far-ultraviolet region (1500-3000 Å) and the near-ultraviolet region only (2000-3000 Å).
}
{
The goodness-of-fit test had a reduced chi-square value of \textasciitilde{}25 for both regions it was performed on. These high values suggest that the model is not representing the observed data well.
}
{
Thus far, the high reduced chi-square indicates that the model poorly fits the data. These high values could be due to an excess in ultraviolet emission, but alternative explanations, such as underestimated uncertainties and model limitations, cannot be ruled out. Further research would help narrow down the source of the high reduced chi-square values and allow for a clearer interpretation of the data in the ultraviolet region. Additionally, higher-quality ultraviolet data and improved models will be beneficial in confirming whether the observed excess ultraviolet emission is evidence of active accretion in the merger remnant.
}

\newpage
